% SOURCE_AUTHORITY: sga5_fr_workpass.tex, lines 15223--15233 (LF-normalized unit).
% SOURCE_SHA256: 791F4EFFC5E02832D5D77ED03518C8156D6F07E4C8238B03545DB93D883FBB28
Antes de demonstrar estes lemas, indiquemos como se aplicam à demonstração da existência dos sistemas $(K'_\nu)$ e $(h'_\nu)$. Raciocina-se por indução sobre $\nu$, supondo já construídos $K'_0,h'_0,\ldots,K'_{\nu-1},h'_{\nu-1}$; aplicam-se os lemas com
\[
A=\mathbb Z/\ell^{\nu+1}\mathbb Z,
\qquad
I=\ell^\nu\mathbb Z/\ell^{\nu+1}\mathbb Z
\]
(de modo que $I^2=0$),
\[
A_0=\mathbb Z/\ell^\nu\mathbb Z.
\]
Tomemos primeiro $L=K_\nu$ e $M=K'_{\nu-1}$ no primeiro lema. Levando em conta a observação 2, obtemos a existência de $K'_\nu=L'$ com o morfismo de transição $K'_\nu\to K'_{\nu-1}$ dado por $w$, de modo que se satisfaça a condição 1 enunciada acima. O isomorfismo $v$ de $D^-(A)$ permite transportar $h_\nu$ para um endomorfismo de $K'$ no sentido de $D^-(A)$. Pode-se supor que este último provém de um endomorfismo de complexos $h''_\nu:K'_\nu\to K'_\nu$. Apliquemos então o segundo lema com $L=K'_\nu$, $\varphi=h''_\nu$ e $\varphi_0=h'_{\nu-1}$; obtemos a existência de $h'_\nu:K'_\nu\to K'_\nu$ de modo que se satisfaça a condição 2); além disso, a condição 3) será satisfeita. Observe-se que, com as notações do lema 1, se $M^i=0$ para um índice $i$, deduz-se $L'^i=0$ graças ao lema de Nakayama; portanto, todos os complexos $K'_\nu$ são efetivamente nulos fora de um mesmo intervalo de graus.
